\chapter{Supervision record}
\label{app:supervision}

Supervision ran roughly fortnightly from June, with a written brief prepared before each meeting. The table lists the meetings and what each one changed in the project. The direction given at the end of July shaped the whole design and is described in Chapter~\ref{ch:method}.

\begin{center}
\small
\begin{tabular}{p{2.6cm}p{10.2cm}}
\toprule
Date (2026) & What was discussed and what followed \\
\midrule
18 June & First meeting. Get to know the dataset, start reading around text-to-image generation, get the School template. \\
9 July & The two strands of the project agreed: emotion from the text, and analysis of the drawings. A six-part structure for the write-up (abstract, introduction, literature review, methodology, results, discussion), and the instruction to argue from cited evidence and not from opinion. Dataset findings presented alongside: the three filename schemes, the timestamps recovered from the image files, the thirteen pre-festival records, the repeated prompts; questions raised about the original prompt sent to the generator and about names in the data. \\
31 July & The methodological direction, and the timeline for the remaining work: measure how similar the two readings are instead of scoring matches; describe the drawing first and extract the emotion from the description; split faces from scenes and emotional from non-emotional prompts; redact names and keep illness content; build in an evaluation. \\
20 August & The demonstrator shown running live. Timing of a human-rated study discussed; agreed to write up the computational results and treat the rated study as the next step. \\
4 September & Project presentation, with a live demonstration. \\
\bottomrule
\end{tabular}
\end{center}
